% from ega4-16-fr.tex
\oldpage[IV-4]{27}
cas particulier d'un formalisme différentiel pour des espaces annelés munis
d'un faisceau d'algèbres sur le faisceau structural. Cependant, nous n'avons
pas voulu partir de ce point de vue, moins intuitif et moins commode pour les
applications. Il semble d'ailleurs que, pour les diverses espèces de
«variétés», la construction «globale» des $\mathscr{P}^n$ analogue à celle que
nous utilisons ici soit également mieux adaptée aux applications.
\end{enumerate}
\end{remarks}

\subsection{Faisceaux et fibrés tangents relatifs; dérivations.}
\EGAReaderTarget{ega4.c04.subsec.16-5}{IV.16.5}

\begin{env}[16.5.1]
\EGAReaderTarget{ega4.c04.statement.16-5-1}{IV.16.5.1}
Soit $f=(\psi,\theta):X\to S$ un morphisme d'espaces annelés. Pour tout
$\mathcal{O}_X$-Module $\mathscr{F}$, on appelle \emph{$S$-dérivation}
(ou \emph{$(X/S)$-dérivation}, ou \emph{$f$-dérivation}) \emph{de
$\mathcal{O}_X$ dans $\mathscr{F}$} tout homomorphisme de \emph{faisceaux de
groupes additifs} $D:\mathcal{O}_X\to\mathscr{F}$, vérifiant les conditions
suivantes:
\begin{enumerate}
\item[(a)] pour tout ouvert $V$ de $X$, et tout couple de sections
$(t_1,t_2)$ de $\mathcal{O}_X$ au-dessus de $V$, on a
\[
  D(t_1t_2)=t_1D(t_2)+D(t_1)t_2;
  \tag{16.5.1.1}
  \label{IV.16.5.1.1}
\]
\item[(b)] pour tout ouvert $V$ de $X$, toute section $t$ de $\mathcal{O}_X$
au-dessus de $V$ et toute section $s$ de $\mathcal{O}_S$ au-dessus d'un ouvert
$U$ de $S$ tel que $V\subset f^{-1}(U)$, on a
\[
  D((s|V)t)=(s|V)D(t).
  \tag{16.5.1.2}
  \label{IV.16.5.1.2}
\]
\end{enumerate}
Il est clair qu'il revient au même de dire que pour tout $x\in X$,
l'homomorphisme de groupes additifs
$D_x:\mathcal{O}_x\to\mathscr{F}_x$ est une
$\mathcal{O}_{f(x)}$-dérivation.

Une autre interprétation consiste à considérer la $\mathcal{O}_X$-Algèbre
$\mathscr{D}_{\mathcal{O}_X}(\mathscr{F})$ égale à
$\mathcal{O}_X\oplus\mathscr{F}$, la structure d'Algèbre étant définie par la
condition que pour tout ouvert $V$ de $X$, le produit de deux sections de
$\mathcal{O}_X$ (resp. d'une section de $\mathcal{O}_X$ et d'une section de
$\mathscr{F}$) au-dessus de $V$ est défini par la structure d'anneau de
$\Gamma(V,\mathcal{O}_X)$ (resp. la structure de
$\Gamma(V,\mathcal{O}_X)$-module sur $\Gamma(V,\mathscr{F})$), et le produit de
deux sections de $\mathscr{F}$ au-dessus de $V$ est pris égal à 0; alors,
$\mathscr{F}$ est un Idéal de
$\mathscr{D}_{\mathcal{O}_X}(\mathscr{F})$, noyau de l'augmentation canonique
$\mathscr{D}_{\mathcal{O}_X}(\mathscr{F})\to\mathcal{O}_X$, et dire que $D$
est une $S$-dérivation de $\mathcal{O}_X$ dans $\mathscr{F}$ signifie que
$1_{\mathcal{O}_X}+D$ est un $\mathcal{O}_S$-homomorphisme d'Algèbres de
$\mathcal{O}_X$ dans $\mathscr{D}_{\mathcal{O}_X}(\mathscr{F})$, qui, composé
avec l'augmentation, donne $1_{\mathcal{O}_X}$.

Les $S$-dérivations de $\mathcal{O}_X$ dans $\mathscr{F}$ forment évidemment
un $\Gamma(X,\mathcal{O}_X)$-module
$\operatorname{D\acute{e}r}_S(\mathcal{O}_X,\mathscr{F})$.

Lorsque $\mathscr{F}=\mathcal{O}_X$, une $S$-dérivation de $\mathcal{O}_X$ dans
lui-même est simplement appelée une \emph{$S$-dérivation de $\mathcal{O}_X$}.
\end{env}

\begin{proposition}[16.5.2]
\label{IV.16.5.2}
Soient $A$ un anneau, $B$ une $A$-algèbre, $L$ un $B$-module; on pose
$S=\operatorname{Spec}(A)$, $X=\operatorname{Spec}(B)$,
$\mathscr{F}=\widetilde{L}$. Alors l'application $D\mapsto\Gamma(D)$ qui, à
toute $S$-dérivation $D$ de $\mathcal{O}_X$ dans $\mathscr{F}$ fait
correspondre l'application $\Gamma(D):t\mapsto D(t)$ de $B$ dans $L$, est un
isomorphisme de $B$-modules de
$\operatorname{D\acute{e}r}_S(\mathcal{O}_X,\mathscr{F})$ sur
$\operatorname{D\acute{e}r}_A(B,L)$ (cf. (0, 20.1.2)).
\end{proposition}

Cela résulte aussitôt de l'interprétation donnée ci-dessus des $S$-dérivations
en termes
